\subsection*{Testing transfer to larger molecules}

To test accuracy beyond the endpoint's training-size range, we held out the largest 15\% of
molecules within each endpoint source and refitted eight of the nine fixed endpoint
configurations on the remaining 1,160 labels. The MoLFormer--ExtraTrees
combination was not included in this diagnostic. The 205 held-out molecules contain
8--37 heavy atoms; 188 exceed the largest training molecule, which has 17.
The overlapping lower end results from applying the rule within each source.
These labels were available during earlier random-fold development, so this is
an exploratory size diagnostic, not an untouched prospective test. The response
surrogates remain fixed: this comparison tests endpoint transfer conditional on
their predictions, not size extrapolation with all response-source exposure removed.

TabPFN reduces the size-held-out MAE from 1.772 for ExtraTrees to 1.318
kcal mol$^{-1}$ (paired difference $-0.453$; 95\% interval, $-0.648$ to $-0.281$;
Fig.~\ref{fig:size}a). Here its response descriptors also help: without them the
error is 1.959 (response-minus-control difference $-0.640$; interval, $-0.847$ to
$-0.451$). The refits use separately locked input recipes, so the cross-endpoint
contrast does not isolate architecture from representation or preprocessing.
TabM, the linear-plus-neural model and the dual-branch model also reduce
error relative to ExtraTrees, with paired intervals below zero. These results
identify useful alternatives for this size split; they do not imply the same ordering
for every new chemical family (Supplementary Table~15; Supplementary Data~6,
\nolinkurl{size_diagnostic_comparisons.csv}).

Capped glycine and alanine series ask a complementary question about prediction
shape. ExtraTrees approaches plateaus near $-23$ and $-25$ kcal mol$^{-1}$, whereas
several neural endpoints and TabPFN continue to change over the tested range
(Fig.~\ref{fig:size}b,c). For Ac-(Gly)$_n$-NHMe, TabPFN predicts $-40.802$ at
$n=6$ and $-67.194$ at $n=12$, compared with $-23.005$ and $-23.440$ for
ExtraTrees. No experimental peptide values are available here: these are predictions,
not errors, and a more negative or more linear curve is not evidence of greater
physical accuracy. The finite outer bound on the mean prediction is approximately
$[-32.595,5.757]$ for the released trees and $[-658.495,649.922]$ for the fitted
TabPFN ensemble. The latter is much wider but still finite; neither architecture
guarantees asymptotic extensivity.

\begin{figure}[!htbp]
\centering
\includegraphics[width=\textwidth]{F8_size_and_stress.pdf}
\caption{\textbf{Labelled size transfer and unlabelled series answer different questions.}
\textbf{a}, MAE on the 205-molecule size holdout after refitting on the same 1,160
training labels. Filled and open marks denote models with and without responses.
All predictions average three seeds. This diagnostic reuses development labels
under a different split. \textbf{b,c}, Predictions along capped glycine and alanine
series, using the full-data models with responses. Neural and TabPFN curves
average five seeds; ExtraTrees averages three. These curves have no peptide
reference values and therefore do not rank physical accuracy. The abscissa is the
number of residues between the acetyl and N-methylamide caps; $n=0$ is
N-methylacetamide. Colors identify the same families in both panels.}
\label{fig:size}
\end{figure}

Linear alkanes expose an additional limitation. TabPFN remains non-monotonic between
C10 and C14 (3.348 and 3.046 kcal mol$^{-1}$), despite its changed peptide behavior.
Moreover, 27 of the 30 tested alkanes occur in endpoint training, so this series is
not an independent extrapolation benchmark. A source audit confirmed that the
low C14 and C16 training values were copied faithfully from their source; removing
these two labels raises the corresponding tree predictions. We retain the original
labels and report the sensitivity rather than substitute an assumed linear trend
(Supplementary Methods~S20).

Finally, the native TabPFN predictive distributions do not yet provide calibrated
reliability. Equal mixtures of the five fitted distributions give nominal 80\%
and 95\% intervals with only 58.2\% and 80.0\% empirical coverage on External-220;
the corresponding Strict-97 coverages are 55.7\% and 76.3\%. These are predictive
intervals, not the spread across seeds. No calibration or model reselection was
performed on these outcomes (Supplementary Fig.~11 and Table~16).
